\input{../../../../templates/es/preambulo.tex}
\setbeamertemplate{navigation symbols}{}
\setbeamertemplate{footline}[frame number]
\lstset{basicstyle=\ttfamily\fontsize{7.3pt}{8.7pt}\selectfont,breaklines=true,columns=fullflexible,keepspaces=true,aboveskip=4pt,belowskip=0pt}
\date{}
\begin{document}
\begin{frame}[fragile]{Tu primer algoritmo genético}
\centering\Large Extenso: estudio
\par\vspace{1em}\normalsize Prof. Juliho Castillo\par Tec de Monterrey\par
\end{frame}

\begin{frame}[fragile]{Objetivos de aprendizaje}
Representar un candidato, su aptitud y una población. Explicar selección, cruza y mutación. Integrar una actualización generacional. Interpretar corridas con semilla sin atribuirles una garantía.\\[0.8em]
Requisitos: variables, condicionales, ciclos y funciones de Python. Aquí se introducen objetos, comprensiones, arreglos y emparejamiento.
\end{frame}

\begin{frame}[fragile]{Modelo y supuestos}
\[\max_{x\in[-10,10]} f(x),\qquad f(x)=\sin x-0.2|x|.\]
Es un modelo sintético y adimensional de enseñanza, no mediciones de Planta Física. El seno recibe radianes. La evaluación es determinista; mayor aptitud es preferible. El intervalo es una restricción obligatoria.\\[0.7em]
Los candidatos y las perturbaciones se generan para aprender el algoritmo. No se requieren archivos de datos externos.
\end{frame}

\begin{frame}[fragile]{1. Definir la aptitud y su referencia}
NumPy evalúa la función elemento a elemento. Una malla de 400 puntos da una referencia muestreada; el algoritmo no busca consultando esa malla.

\begin{lstlisting}
def objetivo(x: float | np.ndarray) -> float | np.ndarray:
    return np.sin(x) - 0.2 * abs(x)
\end{lstlisting}
\end{frame}

\begin{frame}[fragile]{Resultado 1}
\centering\includegraphics[height=0.58\textheight,width=\textwidth,keepaspectratio]{../../figuras/primer_ejemplo_01_el_paisaje.png}
\par\vspace{0.3em}\raggedright\small Mejor punto de malla: x = +1.378, aptitud = +0.706. Tres cimas interiores y una colina parcial en el extremo izquierdo muestran varias cuencas. El ascenso local puede omitir la cima más alta.
\end{frame}

\begin{frame}[fragile]{Una referencia analítica para estudiar}
El algoritmo no usa derivadas. Para auditar este ejemplo, si $x>0$,
$f'(x)=\cos x-0.2$. En $(0,\pi)$ cambia de positiva a negativa en
$x^*=\arccos(0.2)\approx1.3694384$.
Para $x\in[-\pi,0]$, $f(x)\le0$; para $|x|\ge\pi$,
$f(x)\le1-0.2\pi\approx0.3717<f(x^*)\approx0.7059082$.
Por tanto, $x^*$ es el máximo global en el intervalo declarado. La malla y el
algoritmo genético lo aproximan; ninguno calcula esta demostración.
El candidato final de la semilla 16 está \textbf{cerca} de una cima local, no es su
maximizador exacto.

\textbf{Comprobación:} explica por qué redondear ambas aptitudes a \texttt{0.706} no vuelve
idénticos el punto de la malla de 400 valores y el candidato de la semilla 52.
\end{frame}

\begin{frame}[fragile]{2. Guardar un candidato y su aptitud}
La construcción guarda un cromosoma de un gen y calcula su aptitud. La propiedad lee el gen; la representación imprime signo y tres decimales.

\begin{lstlisting}
class Individuo:

    def __init__(self, lista_genes: list[float]) -> None:
        self.lista_genes = lista_genes
        self.aptitud = objetivo(lista_genes[0])

    @property
    def gen(self) -> float:
        return self.lista_genes[0]

    def __repr__(self) -> str:
        return f'x={self.gen:+.3f} f={self.aptitud:+.3f}'
\end{lstlisting}
\end{frame}

\begin{frame}[fragile]{Crear la población inicial}
El muestreo uniforme crea diez candidatos evaluados. Reinicia la semilla antes de la demostración. Una comprensión repite una expresión; max usa su función key para comparar aptitud.

\begin{lstlisting}
def crear_aleatorio() -> Individuo:
    return Individuo([random.uniform(GEN_MIN, GEN_MAX)])

random.seed(SEMILLA)
poblacion = [crear_aleatorio() for _ in range(TAMANO_POBLACION)]
\end{lstlisting}
\end{frame}

\begin{frame}[fragile]{Resultado 2}
\centering\includegraphics[height=0.58\textheight,width=\textwidth,keepaspectratio]{../../figuras/primer_ejemplo_02_poblacion_aleatoria.png}
\par\vspace{0.3em}\raggedright\small La semilla 52 da como mejor inicial x = -0.323, aptitud = -0.382. Diez candidatos aleatorios ofrecen alternativas, sin garantizar cercanía a la mejor cima.
\end{frame}

\begin{frame}[fragile]{Por qué un torneo favorece a los mejores candidatos (1)}
Con $n=10$ candidatos y $k=3$ sorteos independientes \textbf{con reemplazo}, ordena las
aptitudes distintas por rango desde $r=1$ (peor) hasta $r=n$ (mejor).
La probabilidad de que todos tengan rango a lo sumo $r$ es $(r/n)^k$.
Resta el evento de que todos tengan rango a lo sumo $r-1$:
\[P(\text{gana el rango }r)=(r/n)^k-((r-1)/n)^k.\]
El mejor gana un lugar con probabilidad $1-0.9^3=0.271$; su conteo esperado en diez
lugares es $2.71$, no una cantidad garantizada en esta semilla. Los empates requieren
analizar su resolución; la fórmula supone rangos distintos.
\end{frame}

\begin{frame}[fragile]{Por qué un torneo favorece a los mejores candidatos (2)}
Antes de leer la sintaxis compacta: \texttt{[v*v for v in [2,3]]} produce \texttt{[4,9]}.
\texttt{max(["a", "larga"], key=len)} devuelve \texttt{"larga"} al comparar longitudes.
La selección aplica estas ideas a objetos sorteados y sus aptitudes.
\end{frame}

\begin{frame}[fragile]{3. Selección por torneo}
Sortea tres candidatos con reemplazo y elige la mayor aptitud. Repite por cada lugar de salida. El resultado contiene referencias a objetos existentes; no crea genes nuevos.

\begin{lstlisting}
def seleccion_torneo(poblacion: list[Individuo], tamano: int) -> list[Individuo]:
    descendencia = []
    for _ in range(len(poblacion)):
        candidatos = [random.choice(poblacion) for _ in range(tamano)]
        descendencia.append(max(candidatos, key=lambda ind: ind.aptitud))
    return descendencia
\end{lstlisting}
\end{frame}

\begin{frame}[fragile]{Resultado 3}
\centering\includegraphics[height=0.58\textheight,width=\textwidth,keepaspectratio]{../../figuras/primer_ejemplo_03_seleccion.png}
\par\vspace{0.3em}\raggedright\small Cuatro de diez individuos originales quedan con cero copias. Los genes seleccionados son un subconjunto de los originales. Cambian las multiplicidades, no los valores de los genes.
\end{frame}

\begin{frame}[fragile]{Cruza por mezcla y los dos intervalos}
Sean $a,b$ los genes parentales, $d=|b-a|$ y $\alpha\ge0$.
Sortea $u\sim U[0,1)$ y define $s=(1+2\alpha)u-\alpha$.
Las propuestas son $c_1=(1-s)a+sb$ y $c_2=sa+(1-s)b$.
Cada propuesta marginal recorre $[\min(a,b)-\alpha d,\max(a,b)+\alpha d]$.
Comparten un sorteo: son \textbf{complementarias, no independientes} y cumplen
$c_1+c_2=a+b$ antes de acotar. Para $a=-2,b=3,\alpha=1$, el intervalo es $[-7,8]$.

La proyección $\min(10,\max(-10,c))$ impone el dominio legal y puede romper esa
identidad de suma. Si los progenitores son iguales, $d=0$ y esta cruza no explora.
\texttt{par = (2,3)} contiene dos valores; \texttt{izquierdo,derecho = par} los desempaqueta.
\texttt{elementos=[]; elementos.extend(par)} forma \texttt{[2,3]}; \texttt{append(par)} guarda un solo par.
\end{frame}

\begin{frame}[fragile]{Proyectar una propuesta al intervalo legal}
Una propuesta mayor a 10 queda en 10; una menor a -10 queda en -10. Un valor legal no cambia. La proyección forma parte del algoritmo, no es un ajuste visual.

\begin{lstlisting}
def acotar(gen: float, minimo: float=GEN_MIN, maximo: float=GEN_MAX) -> float:
    return max(minimo, min(maximo, gen))
\end{lstlisting}
\end{frame}

\begin{frame}[fragile]{4. Cruza por mezcla}
Sortea $s=(1+2\alpha)u-\alpha$, $u\sim U[0,1)$. Propone $(1-s)a+sb$ y $sa+(1-s)b$; acota al dominio legal. Un sorteo compartido hace propuestas complementarias, no independientes.

\begin{lstlisting}
def cruza_mezcla(gen1: float, gen2: float, alfa: float) -> tuple[float, float]:
    desplazamiento = (1 + 2 * alfa) * random.random() - alfa
    hijo1 = (1 - desplazamiento) * gen1 + desplazamiento * gen2
    hijo2 = desplazamiento * gen1 + (1 - desplazamiento) * gen2
    return (acotar(hijo1), acotar(hijo2))

def cruza(progenitor1: Individuo, progenitor2: Individuo) -> tuple[Individuo, Individuo]:
    g1, g2 = cruza_mezcla(progenitor1.gen, progenitor2.gen, ALFA_MEZCLA)
    return (Individuo([g1]), Individuo([g2]))
\end{lstlisting}
\end{frame}

\begin{frame}[fragile]{Resultado 4}
\centering\includegraphics[height=0.58\textheight,width=\textwidth,keepaspectratio]{../../figuras/primer_ejemplo_04_cruza.png}
\par\vspace{0.3em}\raggedright\small Con progenitores -2 y +3 y alfa = 1, 14 de 20 hijos salen del intervalo parental. Todos permanecen en el dominio legal. Crear valores no garantiza mejorar la aptitud.
\end{frame}

\begin{frame}[fragile]{Mutación gaussiana, límites y conteo observado (1)}
Sortea $\varepsilon\sim N(\mu,\sigma^2)$ y propone
$x'=\min(10,\max(-10,x+\varepsilon))$.
Aquí $\mu=0$ y $\sigma$ es una desviación estándar, no una distancia máxima.
La propuesta sin acotar tiene densidad
\[p(y)=\frac{1}{\sigma\sqrt{2\pi}}\exp\!\left[-\frac{(y-x-\mu)^2}{2\sigma^2}\right].\]
Acotar manda las propuestas fuera del dominio a los extremos; la distribución
resultante ya no es una gaussiana ordinaria. La curva del script describe la
\textbf{propuesta sin acotar}; el histograma del notebook muestra valores muestreados
y acotados. Ambos usan las mismas propuestas sorteadas y el mismo criterio de llegada.
\end{frame}

\begin{frame}[fragile]{Mutación gaussiana, límites y conteo observado (2)}
El experimento cuenta $|x'-h|<1.5$, con centro fijo $h=1.38$,
partiendo de $x=-4.6$. \texttt{110/2000} es la fracción observada \texttt{0.055}, no una probabilidad
establecida para cualquier corrida. Es un experimento de un salto; el algoritmo
completo también decide si muta con probabilidad \texttt{0.1}.
\end{frame}

\begin{frame}[fragile]{5. Mutación gaussiana}
Suma $\varepsilon\sim N(0,\sigma^2)$ y acota. Sigma determina el salto típico; no es una distancia máxima. En el algoritmo completo, la mutación se activa por separado con probabilidad 0.1.

\begin{lstlisting}
def mutacion_gaussiana(gen: float, mu: float, sigma: float) -> float:
    return acotar(gen + random.gauss(mu, sigma))

def mutar(individuo: Individuo) -> Individuo:
    return Individuo([mutacion_gaussiana(individuo.gen, MUTACION_MU, MUTACION_SIGMA)])
\end{lstlisting}
\end{frame}

\begin{frame}[fragile]{Densidad de la propuesta sin acotar}
La curva es la densidad antes de acotar; el recorte agrega masa en los extremos. El notebook muestra histogramas de propuestas muestreadas y acotadas. Ambos usan el mismo criterio de llegada.

\begin{lstlisting}
def densidad_gaussiana(malla: np.ndarray, mu: float, sigma: float) -> np.ndarray:
    return np.exp(-0.5 * ((malla - mu) / sigma) ** 2) / (sigma * np.sqrt(2.0 * np.pi))
\end{lstlisting}
\end{frame}

\begin{frame}[fragile]{Resultado 5}
\centering\includegraphics[height=0.58\textheight,width=\textwidth,keepaspectratio]{../../figuras/primer_ejemplo_05_mutacion.png}
\par\vspace{0.3em}\raggedright\small Desde x = -4.6, cuenta llegadas con |x - 1.38| < 1.5 en 2000 intentos. Sigma 1 da 0; sigma 3 da 110 (5.5\%). Cero eventos observados no significa probabilidad cero.
\end{frame}

\begin{frame}[fragile]{Actualización completa e invariantes (1)}
La población tiene diez miembros. Cada generación ejecuta diez torneos de tamaño
tres, luego cinco decisiones de cruza con probabilidad $0.8$ y diez decisiones de
mutación con probabilidad $0.1$. Reemplaza toda la población. Diez generaciones son
un presupuesto fijo, no una prueba de convergencia. Se esperan cuatro pares cruzados
y un individuo mutado por generación; los conteos efectivos varían.
\end{frame}

\begin{frame}[fragile]{Actualización completa e invariantes (2)}
Si \texttt{lugares=[0,1,2,3]}, \texttt{lugares[::2]} da \texttt{[0,2]}, \texttt{lugares[1::2]} da \texttt{[1,3]} y
\texttt{list(zip(lugares[::2],lugares[1::2]))} da \texttt{[(0,1),(2,3)]}.
Este emparejamiento requiere un tamaño par. Se conservan tamaño de población,
límites y aptitud almacenada. Si un operador no se aplica, puede reutilizarse un
objeto: los operadores no deben mutar en sitio progenitores compartidos.
No hay elitismo; la mejor aptitud de la generación puede disminuir.
\end{frame}

\begin{frame}[fragile]{6. Actualización generacional}
Inicializa diez candidatos; repite selección, cruza de pares, mutación y reemplazo diez veces. La cruza tiene probabilidad 0.8 por par y la mutación 0.1 por individuo. No hay elitismo que proteja al mejor previo.

\begin{lstlisting}
random.seed(SEMILLA)
poblacion = [crear_aleatorio() for _ in range(TAMANO_POBLACION)]
historia = [max(poblacion, key=lambda i: i.aptitud).aptitud]
for generacion in range(1, MAX_GENERACIONES + 1):
    descendencia = seleccion_torneo(poblacion, TAMANO_TORNEO)
    cruzados = []
    for p1, p2 in zip(descendencia[::2], descendencia[1::2]):
        if random.random() < PROB_CRUZA:
            cruzados.extend(cruza(p1, p2))
        else:
            cruzados.extend([p1, p2])
    poblacion = [mutar(i) if random.random() < PROB_MUTACION else i for i in cruzados]
    mejor = max(poblacion, key=lambda i: i.aptitud)
    historia.append(mejor.aptitud)
    print(f'Generación {generacion:2d} | mejor {mejor.aptitud:+.4f}')
\end{lstlisting}
\end{frame}

\begin{frame}[fragile]{Resultado 6}
\centering\includegraphics[height=0.58\textheight,width=\textwidth,keepaspectratio]{../../figuras/primer_ejemplo_06_el_ciclo_completo.png}
\par\vspace{0.3em}\raggedright\small Semilla 52: x final = +1.372, aptitud = +0.706. La malla reporta x = +1.378. Son aproximaciones diferentes aunque sus aptitudes mostradas se redondeen igual.
\end{frame}

\begin{frame}[fragile]{7. Cambiar sólo la semilla}
Cambiar 52 por 16 cambia la inicialización y cada decisión aleatoria posterior. Conserva operadores, tamaño, probabilidades y diez generaciones. No se aísla la inicialización como causa única.
\end{frame}

\begin{frame}[fragile]{Resultado 7}
\centering\includegraphics[height=0.58\textheight,width=\textwidth,keepaspectratio]{../../figuras/primer_ejemplo_07_optimo_local.png}
\par\vspace{0.3em}\raggedright\small Semilla 16: x final = -4.417, aptitud = +0.073, cerca de una cima local menor. Esta corrida de diez generaciones omite la cima más alta; no demuestra que nunca pueda escapar.
\end{frame}

\begin{frame}[fragile]{Comprobaciones e interpretación}
Comprueba tamaño de población, límites de genes y concordancia de cada gen con su aptitud almacenada. Repetir una semilla debe reproducir la población ordenada y el historial completo.\\[0.7em]
Los conteos son exactos. Entre entornos, compara valores sin redondear con $\mathrm{rtol}=10^{-12}$ y $\mathrm{atol}=10^{-12}$. Investiga trayectorias distintas. Dos semillas no estiman una tasa de éxito.
\end{frame}

\begin{frame}[fragile]{Práctica: predecir, modificar, explicar}
1. Conserva los progenitores y compara $\alpha=0,0.5,1$. ¿Qué cambia el intervalo de propuestas?\\[0.6em]
2. Prueba $\sigma=2$ desde el mismo origen. ¿El conteo determina la mejor configuración de una corrida completa?\\[0.6em]
3. Ejecuta las semillas 0, 1, 2 y 3. Define éxito antes de diseñar un estudio mayor. Usa nombres separados y conserva la línea base.
\end{frame}

\begin{frame}[fragile]{Respuestas razonadas}
1. Aumentar alfa amplía el intervalo. No asegura mejor aptitud: un hijo puede caer en un valle.\\[0.6em]
2. Sigma cambia el alcance de un salto. El desempeño completo también depende de selección, cruza, activación de mutación y presupuesto generacional.\\[0.6em]
3. Cuatro corridas ilustran variación. Estimar una tasa requiere un evento explícito, suficientes corridas independientes e incertidumbre reportada.
\end{frame}

\begin{frame}[fragile]{Conexión con Planta Física}
Un candidato codificará decisiones de diseño físico; la factibilidad reflejará restricciones de ingeniería; la aptitud medirá un criterio de desempeño justificado.\\[0.7em]
La función de un gen es un modelo didáctico. Su óptimo numérico no es una recomendación para una instalación real. Distingue el mejor diseño observado de un óptimo global verificado.
\end{frame}

\begin{frame}[fragile]{Arreglos y gráficas}
\texttt{linspace} crea valores alineados; \texttt{argmax} devuelve un índice. Figure es el lienzo completo; Axes es un área de dibujo. Guarda, muestra cuando corresponda y cierra la figura. Los rótulos indican qué mide cada eje.
\end{frame}

\begin{frame}[fragile]{Referencias y conteos}
La selección devuelve referencias compartidas. La identidad difiere de la igualdad de genes. Counter cuenta identificadores; get(key, 0) contempla individuos ausentes. La variación crea objetos nuevos para conservar la aptitud correcta.
\end{frame}

\begin{frame}[fragile]{Pares y reemplazo}
\texttt{[::2]} selecciona lugares 0,2,...; \texttt{[1::2]} selecciona 1,3,... . zip empareja posiciones. extend agrega elementos; append agrega un objeto. Una comprensión condicional muta o conserva el objeto; la asignación reemplaza la lista de población.
\end{frame}

\begin{frame}[fragile]{APIs de gráficas y escala}
Alfa de dibujo controla opacidad, no alcance de cruza. twinx comparte x con otra escala y para densidad; combina leyendas explícitamente. Los histogramas cuentan observaciones por intervalo. Una altura de densidad no es probabilidad. Acompaña colores con estilos y rótulos.
\end{frame}

\begin{frame}[fragile]{Tipos, estado y reproducibilidad}
Un tipo anotado documenta entradas y salidas, sin imponerlas durante la ejecución. self accede al estado del objeto; una propiedad expone un atributo calculado. La semilla fija los sorteos posteriores: cambiar su orden cambia la corrida. Reinicia kernels antes de validar.
\end{frame}

\begin{frame}[fragile]{Fuente y profundización}
Ivan Gridin, \emph{Learning Genetic Algorithms with Python}, capítulo 1, aporta la base. El curso desarrolla siete incrementos autónomos.\\[0.7em]
El notebook extenso contiene el recorrido ejecutable completo, ejemplos de sintaxis, experimentos y soluciones. La Lección 02 nombra las fases del mismo ciclo; esta lección conserva su línea base.
\end{frame}
\end{document}
